% ── 全文通用形态（Mrite 高教社杯规范）──
% 正文一律自然段；禁分点符号、禁 \textbf（摘要“针对问题X”与问题重述“问题N：”除外）。
% 跨页表用 longtable，短表用 tabularx；列宽比例总和 = 1.04 − 0.04 × N（N 为列数），格内文字不换行。
% 图宽 0.8\textwidth，表宽 \textwidth；图题交给 \caption，绘图代码不写 set_title()；图表必须在正文被引用。
% 章与章、节与节之间一律不加 \newpage。
% ── 本节合同 ──
% 总字数 ≤900，开头段 ≤30 字，结尾段 ≤20 字；编译后必须恰为 1 页（\label{abstract:end} 供 .aux 程序化核验，不得删）。
% 每问一段，段内五动作固定：段首定位 → 建模主张 → 方法链条 → 结果句 → 验证句。
% 加粗只允许三类、每类至多一处：建立**XX模型** / 用**XX算法** / 得出**XXX数字或结论**。其余一律不加粗。
% 详细 move 库与 12 篇国奖语料出处见 references/abstract-moves.md（写摘要前读）。
%
\begin{abstract}

% ==== 开头段（2句话）====
...是...领域的重要研究课题。本文基于...数据，围绕...展开系统研究。

% ==== 逐问摘要 ====
\textbf{对于问题一,} 首先对...进行了...分析，接着采用...方法构建了...模型，最后通过...得到...。结果表明，...（关键数值）。

\textbf{对于问题二,} 首先...，接着...，最后...。结果表明...。

\textbf{对于问题三,} 首先...，接着...，最后...。结果表明...。

\textbf{对于问题四,} 首先...，接着...，最后...。结果表明...。

% ==== 结尾段（2句话）====
本文完成了...的研究，各模型性能优越、结论可靠。研究成果可为...提供理论依据和决策参考。


\keywords{关键词1；关键词2；关键词3；关键词4；关键词5}

\end{abstract}
